\section{Uniform frequency bounds without locating the zeros}

This section treats the large-frequency issue by discarding a finite
input packet containing all nearby zeros, on both sides of the line.
The discarded packet has logarithmic rank.  Its purpose here is only
an upper bound; no claim that a multiplier acts by a scalar on an
off-line packet is needed.

Put $F_\chi(z)=L(1/2+iz,\chi)$ and $H=2+|T|$.
A zero $\rho=\beta+i\gamma$ corresponds to
$z_\rho=\gamma-i(\beta-1/2)$.
Thus, at large ordinates, every zero lies in $|\Im z|<1/2$.
All constants below may depend on the fixed character and on any
fixed derivative budget, but not on distances between zeros or
their distances to the real axis.

\begin{lemma}\label{lem:rectangle-parametrix}
For sufficiently large $|T|$ there is a real function
$a_T\in C_c^\infty((T-2,T+2))$, with $0\leq a_T\leq1$ and
$a_T=1$ on $[T-1/2,T+1/2]$, and a packet of
$N\leq C_\chi\log H$ zeros $z_1,\ldots,z_N$, counted with
multiplicity, having the following properties.
Every zero whose real projection belongs to the support of $a_T$
belongs to the packet, and $a_T$ is identically one near that
projection.  For each fixed integer $j\geq0$, let $R_{T,j}$ be
the sum of the principal parts at these zeros of
$z^j/F_\chi(z)$.  Give real poles the negative-time boundary
prescription.  Then
\[
 r_{T,j}(\tau)=\frac{\tau^j a_T(\tau)}{F_\chi(\tau)}
                                      -R_{T,j}(\tau)
\]
is an ordinary smooth function on the real line and, for every
fixed $m\geq0$,
\begin{equation}\label{eq:rectangle-remainder}
                 \|r_{T,j}^{(m)}\|_2\leq C_{j,m}H^{C_{j,m}}.
\end{equation}
For each fixed $A>0$, if $R_{T,j,+}$ and $R_{T,j,-}$ denote the
exponential-polynomial inverse transforms on the positive and
negative time half-lines, extended as exponential polynomials to
all real t, then
\begin{equation}\label{eq:rectangle-residue-difference}
 \sup_{-A\leq t\leq0}
 |R_{T,j,-}(t)-R_{T,j,+}(t)|\leq C_{j,A}H^{C_{j,A}}.
\end{equation}
The same assertions hold for a numerator analytic on a fixed
neighborhood of the rectangle below and polynomially bounded in H.
\end{lemma}

\begin{proof}
We first obtain a uniformly bounded reciprocal away from the local
zeros.  On a sufficiently large fixed disk centered at $2+iT$ in
the s-plane, the usual vertical-strip estimate gives a polynomial
upper bound for $L(s,\chi)$; its value at the center is bounded
below by its absolutely convergent Euler product.  Jensen's formula
therefore bounds the number $N_0$ of zeros in a slightly smaller
fixed disk by $C_\chi\log H$.
Divide by the disk Blaschke product of those zeros.  The resulting
factor is zero-free.  The maximum principle gives its polynomial
upper bound, and Harnack's inequality applied to a large constant
minus its logarithmic modulus gives a polynomial reciprocal bound
on a fixed inner disk containing
\[
 \{z:|\Re z-T|\leq2,\quad |\Im z|\leq3/4\}
\]
in the z-coordinate.  All disks may be chosen with fixed radii,
for example outer radius eight, a Blaschke disk radius between six
and seven, and an inner radius five in the s-coordinate.
For $|T|$ large they contain no trivial zero or principal pole.

Let $\gamma_1,\ldots,\gamma_{N_0}$ be the real projections of the
zeros in the Blaschke disk, with multiplicity.  For either interval
\[
 I_-=[T-3/2,T-1],\qquad I_+=[T+1,T+3/2],
\]
the integral of $\log\min(1,|x-\gamma_k|)$ is bounded below by an
absolute constant, uniformly in $\gamma_k$.
Summing and averaging produces points $x_-\in I_-$ and
$x_+\in I_+$ such that
\[
 \prod_{k=1}^{N_0}\min(1,|x_\pm-\gamma_k|)\geq e^{-C N_0}.
\]
Every factor is at most one; consequently each individual distance
is at least $e^{-CN_0}$.  Choose $\varepsilon>0$, bounded below by
$H^{-C_\chi}$, so small that on
$|x-x_\pm|\leq4\varepsilon$ these product bounds remain at least
$e^{-C' N_0}$.  The vertical strips above these intervals are
zero-free and have polynomial reciprocal L bounds, since
$|z-z_k|\geq|\Re z-\gamma_k|$ in every Blaschke factor.

Take the rectangle $\Gamma$ with vertical sides $x_-,x_+$ and
horizontal sides $\Im z=\pm3/4$.  On the horizontal sides every
nontrivial zero is at distance at least $1/4$; the same Blaschke
factorization therefore bounds the reciprocal there polynomially.
Thus $|1/F_\chi(z)|\leq H^C$ on the contour and on suitable
neighborhoods of its sides.  The packet consists of all zeros
inside it.  Choose $a_T$ equal to one on
$[x_-+\varepsilon,x_+-\varepsilon]$, zero outside
$[x_--\varepsilon,x_++\varepsilon]$, with transitions inside the
two zero-free strips.  Standard scaled cutoffs give
$\|a_T^{(m)}\|_\infty\leq C_m\varepsilon^{-m}\leq H^{C_m}$
for each fixed m.  This proves the support and plateau assertions.

Let $R_{T,j}$ be the proper rational function specified in the
statement.  Inside the rectangle,
$z^j/F_\chi(z)-R_{T,j}(z)$ is holomorphic and is represented by
the Cauchy integral on $\Gamma$.  Away from the transition strips,
this integral and its fixed derivatives have polynomial bounds.
Within a transition strip, move the vertical contour to its other
side, without crossing a zero.  The two possible contours are
separated by a distance comparable to $\varepsilon$, so Cauchy's
inequality gives the same polynomial bounds for fixed derivatives
there.  This also bounds $R_{T,j}$ and its derivatives on those
zero-free strips.  Multiplication by $a_T$ now shows that the
remainder is smooth and polynomially bounded on $|\tau-T|\leq3$.
The real pole prescriptions cancel because $a_T=1$ near every real
pole in the packet.

For completeness the bounds at infinity can be obtained without
dividing by a gap between roots.  Write
$P(z)=\prod_{k=1}^N(z-z_k)$.  The numerator of $R_{T,j}$ is a
polynomial of degree less than N. More explicitly, with the contour
oriented counterclockwise,
\[
 P(z)R_{T,j}(z)=\frac1{2\pi i}\int_\Gamma
 \frac{w^j}{F_\chi(w)}\frac{P(z)-P(w)}{z-w}\,dw.
\]
For $z$ outside the contour, the residue formula gives this identity:
the omitted integral of $P(w)w^j/(F_\chi(w)(z-w))$ is zero because
its integrand is holomorphic inside. Both sides are polynomials in
$z$, so the identity holds everywhere. The divided difference
$(P(z)-P(w))/(z-w)$ is a polynomial of degree at most $N-1$ in
$z$; its coefficients in powers of $z-T$ are bounded by $C^N$
on the fixed-size contour. The contour bound therefore bounds the
numerator coefficients by $H^{C_j}C^N$, without an inverse root
distance. Since $N=O(\log H)$ this is a polynomial bound in H.
Outside $|\tau-T|\leq3$, all roots have
$|\Re z_k-T|\leq3/2$ and $|\Im z_k|<1/2$; hence, for each fixed m,
\[
 |R_{T,j}^{(m)}(\tau)|
 \leq H^{C_{j,m}}(1+|\tau-T|)^{-m-1}.
\]
Here factors $C^N$ and fixed powers of N are again absorbed in H.
The smooth remainder equals $-R_{T,j}$ in this region,
which proves \eqref{eq:rectangle-remainder}.

For positive time, closing the inverse-transform contour downward
gives minus i times the residues at the lower-half-plane poles.
For negative time, closing upward gives plus i times the upper
poles, with real poles assigned to that side.  Therefore
\[
 R_{T,j,-}(t)-R_{T,j,+}(t)
       =\frac1{2\pi}\int_\Gamma
                     \frac{z^j e^{-izt}}{F_\chi(z)}\,dz.
\]
On $-A\leq t\leq0$ the exponential has modulus at most
$e^{3A/4}$.  The contour length is bounded, and its integrand has
the polynomial bound already proved.  This establishes
\eqref{eq:rectangle-residue-difference}, including multiple poles.
An analytic numerator obeys the identical contour estimates.
\end{proof}

\begin{lemma}\label{lem:allpole-boundary-bound}
Let $W_{L,T}$ be the orthogonal projection in $S_LL^2(\mathbb R)$
onto functions whose Fourier transforms vanish at every packet
zero $z_k$ to its specified multiplicity.  Equivalently it removes
the span of $1_{[-L,L]}(t)t^j e^{-i\overline z_k t}$.
Set
\[
 b_T=a_T/F_\chi,\qquad T_{L,T}=U^{-1}C_{b_T}W_{L,T}.
\]
For every fixed integer j,
\begin{equation}\label{eq:allpole-bulk-size}
 \|UT_{L,T}\|_{\mathcal S_2(L^2,H^j)}
       \leq H^{C_j}(1+L)^{N+C_j},\qquad L\geq1.
\end{equation}
For every fixed smooth endpoint multiplier $m(t)$, vanishing for
$t<-A$ and with all derivatives bounded, and for every fixed J,j,
the upper-translated boundary operator
\[
 m(t)\,\tau_{-L}UT_{L,T}\tau_L
\]
has Hilbert--Schmidt norm bounded by $H^{C_{J,j,A}}$ after any
output derivative of order at most j and output weight
$(1+|t|)^J$.
The same assertion holds for the raw Fourier-dual endpoint,
with the reflected packet and the fixed conductor translation.
\end{lemma}

\begin{proof}
We first bound the bulk rational inverse.  Every proper rational
function with denominator P and numerator coefficients as in the
previous proof has a Newton expansion
\[
        R_{T,0}(z)=\sum_{j=1}^N
              c_j\prod_{k=1}^j(z-z_k)^{-1},\qquad |c_j|\leq H^C.
\]
The coefficient bound has no gap denominator: the divided
differences of a polynomial are bounded by its coefficients times
binomial coefficients and fixed powers of the bounded roots.
All these factors are $C^N$ and are absorbed in H.

On an input supported in $[-L,L]$ and vanishing at $z_k$, the
inverse of the first-order multiplier $z-z_k$ can be realized by
an exponential integration kernel of modulus at most one, using
the decaying orientation.  The zero moment makes its output
vanish beyond the opposite endpoint as well.  It therefore
remains supported in $[-L,L]$, has operator norm at most $2L$,
and Hilbert--Schmidt norm at most $2L$ on that interval.
Division removes one zero and preserves all remaining zero
conditions, including repeated ones.  Successive division hence
bounds a j-factor term by $(2L)^j$ in Hilbert--Schmidt norm.
This also proves the assertion for real poles, for which the
specified distributional prescription gives the same compact
solution after its zero moment is imposed.
The remainder $r_{T,0}$ is an L² multiplier; convolution by its
inverse transform followed by $W_{L,T}$ has squared
Hilbert--Schmidt norm at most
$2L\|r_{T,0}\|_2^2/(2\pi)$.
This proves the j=0 case of \eqref{eq:allpole-bulk-size}.
The output multiplier $b_T$ is supported in $[T-2,T+2]$, so
fixed logarithmic Sobolev powers cost only an additional fixed
power of H.  Alternatively repeat the same argument with
$R_{T,j}$ and $r_{T,j}$.

For the boundary estimate let $\kappa_T$ be the inverse transform
of $b_T$.  After translating the upper endpoint, the input lies
in $[-2L,0]$ and still satisfies the identical exponential-moment
conditions; translating a polynomial jet preserves its span.
For each output derivative j, decompose the inverse transform
of $\tau^j b_T$ as the inverse transform of $r_{T,j}$ plus the
two exponential polynomials $R_{T,j,+},R_{T,j,-}$ on their
respective half-lines.
The whole exponential polynomial $R_{T,j,+}(t-s)$ annihilates
the input projection exactly.  Subtract it before estimating.
The rational contribution which remains is
\[
 1_{t-s<0}\bigl(R_{T,j,-}(t-s)-R_{T,j,+}(t-s)\bigr).
\]
On the quadrant $t\geq-A$, $s\leq0$, its support is the bounded
triangle $-A\leq t\leq s\leq0$.
The contour bound \eqref{eq:rectangle-residue-difference} therefore
gives a polynomial Hilbert--Schmidt bound, with every stated
fixed output weight and derivative.
This is the step that cancels the large individual residues of
poles approaching each other from opposite sides of the axis.

For the smooth remainder, Plancherel and
\eqref{eq:rectangle-remainder} give every fixed polynomially
weighted L² time norm.  In particular
\[
 \int_{t\geq-A,\ s\leq0}(1+|t|)^{2J}
       |\check r_{T,j}(t-s)|^2\,dt\,ds
 \leq C_{J,A}\int_{u\geq-A}(1+|u|)^{2J+1}
                         |\check r_{T,j}(u)|^2\,du
 \leq H^{C_{J,j,A}}.
\]
Input restriction and right composition by the orthogonal
projection cannot increase this Hilbert--Schmidt norm.
Leibniz's rule includes derivatives of m.
Using the principal part of $\tau^j/F$ at each derivative stage
is essential: polynomial terms at frequency infinity are thereby
absorbed exactly, so no spurious delta derivative at time zero
is omitted.  The equality is first distributional and then an
equality of the locally square-integrable kernels used above.

The raw Fourier transform acts in Mellin coordinates by reflection
and a fixed smooth gamma multiplier.  The functional equation
identifies its inverse multiplier with the dual-character
reciprocal, times an analytic conductor factor.  Reflection of
the input zero conditions gives exactly the dual packet.
The previous lemma's analytic-numerator version therefore proves
the same bounds at that endpoint.  Ordinary raw derivatives of
the scaled boundary insert factors $e^{-t}$; these are bounded
on $t\geq-A$, so their fixed budgets are covered as well.
\end{proof}

\paragraph{A pair approaching the line.}
The elementary example $F(z)=z^2+a^2$, $0<a\leq1/2$, shows why
the packet must contain both sides of the line.  Its inverse
transform is $e^{-a|t|}/(2a)$; each half-line residue separately
diverges as $a\downarrow0$.  On inputs supported in $[-2L,0]$
and orthogonal to $e^{at}$ and $e^{-at}$, subtract the entire
positive-time exponential.  On the upper boundary quadrant the
remaining rational kernel is exactly
\[
          1_{t<s}\frac{\sinh(a(t-s))}{a}.
\]
It is supported on $-A\leq t\leq s\leq0$ and is uniformly
bounded there by $A e^{A/2}$.  It converges to
$1_{t<s}(t-s)$ as the pair reaches the line.  The contour argument
above is the finite-packet version of this cancellation and
retains all multiplicities.

\begin{theorem}\label{thm:unconditional-frequency-tightness}
For the fixed primitive Dirichlet prolate cutoff, including the
two-constraint principal sector, there are constants C,A such that
the positive defect measures satisfy
\[
            \mu_\Lambda([T,T+1])\leq C(2+|T|)^A
            \qquad(\Lambda\geq2,\ T\in\mathbb R).
\]
No hypothesis on zero location or zero separation is used.
Consequently the local convergence of
Theorem~\ref{thm:rational-local-limit} extends to every Schwartz
frequency test.
\end{theorem}

\begin{proof}
Use the preceding bump centered at $T+1/2$ and the full local
zero packet, whose projection we denote by $P_{L,T}$.
The construction proving the local theorem applies with
$W_{L,T}=S_L-P_{L,T}$: the extra off-line zero conditions do not
alter the exact multiplier identity $E_\chi T_{L,T}=C_{a_T}W_{L,T}$.
By Lemma~\ref{lem:allpole-boundary-bound}, all its endpoint
operators, with the finite derivative and moment budgets used in
the cross estimates, have Hilbert--Schmidt norms at most a fixed
power of H, uniformly in L.  Their theta images have the same
property by \eqref{eq:local-theta-graph}.
The outside-input boundary $S_LC_{a_T}(1-S_L)$ has the same bound
by the fixed derivative bounds for $a_T$.

In the principal sector its two additional constraint rows are
bounded by the analytic-numerator version of the same argument.
Here is the continuation needed for that application.
For $k=1,2$, put $\theta_k(t)=e^{t/2}p(e^t)^k$.
The function $(\partial_t-1/2)\theta_k$ is smooth and compactly
supported, because $p(e^t)=1$ near negative infinity.
Consequently
\[
 \widehat\theta_k(z)=
 -\frac{\widehat{(\partial_t-1/2)\theta_k}(z)}{1/2+iz}
\]
continues meromorphically, with its only possible pole at $z=i/2$.
On every fixed-height rectangle centered at a large real T it is
analytic and has polynomial bounds; its compactly supported
numerator in fact has rapid horizontal decay there.
The normalized mean row, up to the already estimated cross error,
is the restriction to $s\leq0$ of
$\check{(b_T\widehat\theta_2(-\cdot))}(-s)$, followed by the
translated input projection.  The reflected positive-time residue
polynomial is annihilated by that projection, leaving the smooth
remainder with the polynomial L² bound of
Lemma~\ref{lem:rectangle-parametrix}.  The value row is the same
calculation with $\theta_1$ and the dual inverse.
Only the unweighted input-row $L^2$ bound is needed here, and right
composition with the orthogonal input projection preserves that bound.
Uniform polynomial input moments are not asserted after projection.
The two output polar profiles have fixed L² norms after endpoint
translation, so both additional boundary operators have the
required polynomial Hilbert--Schmidt bounds.

The strict-margin prolate graph approximation and
\eqref{eq:allpole-bulk-size} bound the bulk error by
\[
             H^C(1+L)^{C\log H+C}e^{-c_0L}
\]
after fixing a sufficiently large derivative budget.  The
cross-boundary errors have a similar bound, or the stronger
$H^C e^{-c_0L}$ bound.  There is a fixed D for which all are at
most one when $L\geq H^D$.
No approximation of the boundary ranges by the theta images is
needed for this estimate: keep their polynomial Hilbert--Schmidt
norms and use contractivity of $1-Q_\Lambda$.
It follows that
\[
 \|(1-Q_\Lambda)S_LC_{a_T}(1-P_{L,T})\|_2\leq H^C.
\]
The complementary input term has squared norm at most
$\operatorname{rank}P_{L,T}\leq C\log H$, because
$\|a_T\|_\infty\leq1$.  The two inputs are orthogonal, so their
squared Hilbert--Schmidt norms add.  Since the bump is one on
$[T,T+1]$, the exact positive-defect distance formula gives the
desired bound above the threshold.
Below the threshold use $d\mu_\Lambda/d\tau\leq L/\pi$.
Bounded heights follow from the already established local
convergence and this same elementary bound for the remaining
bounded cutoff range.
Finally sum the polynomial mass bound over unit intervals against
the rapidly decreasing supremum of a Schwartz test.  The tails
are uniformly summable; local convergence then proves the full
Schwartz convergence.
\end{proof}

